%-----------------------------------------------------------------------------%
\chapter{\babDua}
\label{bab:2}
%-----------------------------------------------------------------------------%

Bab ini membahas tinjauan pustaka yang menjadi landasan konseptual dan metodologis dalam penelitian mengenai prediksi risiko pelanggaran \textit{Service Level Agreement }(SLA) pada operasi logistik \textit{third party logistics} (3PL) dengan karakteristik Multi-SLA. Tinjauan pustaka dilakukan secara sistematis untuk mengidentifikasi, mengevaluasi, dan mensintesis penelitian-penelitian terdahulu yang relevan dengan topik penelitian.

Dengan adanya tinjauan pustaka yang komprehensif, penelitian ini diharapkan memiliki dasar teoritis dan empiris yang kuat dalam mengembangkan model prediktif untuk mengestimasi risiko pelanggaran SLA secara kuantitatif.

%-----------------------------------------------------------------------------%
\section{Metode Pemilihan Pustaka}
\label{sec:2_1}
%-----------------------------------------------------------------------------%

Pada penelitian ini, tinjauan pustaka dilakukan secara sistematis untuk memastikan literatur yang digunakan relevan, mutakhir, dan sesuai dengan tujuan penelitian. Metode yang digunakan adalah \textit{Systematic Literature Review} (SLR) berdasarkan pendekatan Kitchenham dari \citep{kitchenham}, yang banyak digunakan dalam penelitian sistem informasi.

Pendekatan ini dipilih karena memberikan proses yang terstruktur dalam mengidentifikasi, mengevaluasi, dan mensintesis penelitian sebelumnya secara komprehensif. Selain itu, penelitian ini juga mengadopsi prinsip PRISMA (\textit{Preferred Reporting Items for Systematic Reviews and Meta-Analysis}) \citep{prisma2021} untuk meningkatkan transparansi dalam proses seleksi literatur.

Dengan demikian, Kitchenham digunakan sebagai kerangka metodologi utama SLR, sementara PRISMA digunakan sebagai panduan pelaporan proses seleksi literatur secara transparan dan terstruktur.

\subsection{\textit{Systematic Literature Review} (SLR) dengan Pendekatan Kitchenham}

Systematic Literature Review (SLR) merupakan metode penelitian yang bertujuan untuk mengidentifikasi, mengevaluasi, dan menginterpretasikan seluruh penelitian yang relevan terhadap suatu topik secara sistematis. Dalam penelitian ini, SLR digunakan untuk memperoleh pemahaman komprehensif mengenai pendekatan analitik dan prediktif dalam pengelolaan risiko keterlambatan dan SLA pada logistik.

Mengacu pada \citep{kitchenham}, tahapan SLR terdiri dari tiga fase utama: 
\begin{enumerate}
    \item \textbf{Perencanaan (\textit{Planning})} Menentukan tujuan, research question, serta strategi pencarian literatur.
    \item \textbf{Pelaksanaan (\textit{Conducting})} Melakukan pencarian, seleksi berdasarkan judul dan abstrak, serta evaluasi full-text menggunakan kriteria inklusi dan eksklusi.
    \item \textbf{Pelaporan (\textit{Reporting})} Menyajikan hasil kajian secara sistematis, termasuk sintesis penelitian dan identifikasi \textit{research gap}.
\end{enumerate}

Pendekatan ini memastikan bahwa proses kajian literatur dilakukan secara terstruktur dan dapat direplikasi.

\subsection{Strategi Pencarian Pustaka}

Pencarian literatur dilakukan dengan menggunakan beberapa basis data ilmiah yang memiliki reputasi tinggi dan banyak digunakan dalam penelitian teknologi informasi, yaitu ScienceDirect dan Scopus, yang keduanya diakses menggunakan Remote Lib Universitas Indonesia \url{https://remote-lib.ui.ac.id/menu}.

Pemilihan basis data tersebut didasarkan pada cakupan publikasi yang luas, kualitas jurnal yang terindeks, serta relevansi terhadap topik penelitian yang berkaitan dengan logistik, SLA, dan ML. Kata kunci yang digunakan dalam proses pencarian literatur disusun berdasarkan pertanyaan penelitian dan konsep utama yang diteliti. Kata kunci tersebut antara lain: \textit{”service level agreement”, ”SLA Breach”, ”delivery delay”, ”logistics”, ”third-party logistics”, ”predictive analytics”, ”machine learning”, ”risk”} dan \textit{”optimization”}. 

Kata-kata kunci tersebut dikombinasikan menggunakan operator Boolean seperti \textbf{AND} dan \textbf{OR}, juga dengan memperkecil rentang penerbitan hanya dari tahun 2021 hingga 2026, untuk memperoleh hasil pencarian yang lebih spesifik dan relevan. 

\subsection{Kriteria Pencarian}
Untuk memastikan kualitas dan relevansi literatur, penelitian ini menetapkan kriteria inklusi dan eksklusi sebagai berikut:


\begin{enumerate}
    \item Kriteria Inklusi
        \begin{enumerate}
            \item Artikel ilmiah yang dipublikasikan pada tahun 2021-–2026
            \item Terindeks pada basis data Scopus atau ScienceDirect
            \item Membahas topik terkait SLA, kinerja logistik, atau keterlambatan pengiriman
            \item Menggunakan pendekatan analitik, statistik, atau machine learning
        \end{enumerate}
    \item Kriteria Eksklusi
        \begin{enumerate}
            \item Artikel yang tidak berkaitan dengan konteks logistik (misalnya domain kesehatan atau jaringan komunikasi)
            \item Artikel yang tidak membahas kinerja pengiriman atau SLA
            \item Artikel non-ilmiah seperti opini atau laporan tanpa metodologi yang jelas
        \end{enumerate}
\end{enumerate}

\subsection{Tahapan Seleksi Literatur Lanjutan Menggunakan Metode PRISMA}

Setelah proses pencarian literatur dilakukan berdasarkan pendekatan Kitchenham, tahap seleksi literatur dilaksanakan menggunakan kerangka PRISMA untuk memastikan transparansi dan keterlacakan proses seleksi. PRISMA mengelompokkan proses seleksi ke dalam empat tahap utama, yaitu:
\begin{enumerate}
    \item \textit{Identification} \\ Mengidentifikasi seluruh literatur dari basis data yang digunakan.
    \item \textit{Screening} \\ Menyaring literatur berdasarkan judul dan abstrak.
    \item \textit{Eligibility} \\ Mengevaluasi dokumen secara penuh \textit{(full-text review)} berdasarkan kriteria inklusi dan eksklusi. 
    \item \textit{Included} \\ Menentukan literatur akhir yang digunakan dalam penelitian.
\end{enumerate}


Pencarian literatur dilakukan dalam beberapa tingkat, pada tingkat pertama, kata-kata kunci yang digunakan menghasilkan hasil pencarian dalam jumlah besar, sehingga dibutuhkan pencarian dengan kata kunci tingkat kedua yang lebih mendalam.

Sebagai contoh, hasil pencarian pertama pada Scopus, menghasilkan 218 dokumen dengan kata kunci \textit{( "SLA breach" OR "service level violation" OR "delivery delay" ) AND ( "logistics" OR "last mile" OR "e-commerce delivery" ) AND ( "predictive" OR "analytics" OR "machine learning" )} seperti yang terlihat pada Gambar \ref{fig:key1_scopus} di bawah ini.

\begin{figure}
    \centering
    \includegraphics[width=1\linewidth]{MKSK/Tugas 1/assets/pics/keyword1_scopus.png}
    \caption{Pencarian Tingkat 1 di Scopus yang Kurang Spesifik}
    \label{fig:key1_scopus}
\end{figure}

Kemudian, dengan menggunakan kata kunci yang lebih spesifik dan mendalam, ditemukan 45 dokumen hasil pencarian yang lebih mudah untuk disaring dan dipilih sebagai bahan kajian literatur.

\begin{figure}
    \centering
    \includegraphics[width=1\linewidth]{MKSK/Tugas 1/assets/pics/keyword2_scopus.png}
    \caption{Hasil Pencarian Tingkat 2 di Scopus yang Lebih Mendalam}
    \label{fig:key2_scopus}
\end{figure}

Dari kedua basis data yang digunakan, yaitu ScienceDirect dan Scopus, diterapkan proses seleksi menggunakan metode PRISMA yang terdiri dari identifikasi, seleksi atau \textit{screening}, evaluasi keseluruhan dokumen atau \textit{eligibility}, dan pemilihan literatur  \textit{included}.

\begin{figure}
    \centering
    \includegraphics[width=1\linewidth]{MKSK/Tugas 1/assets/pics/prisma.png}
    \caption{Tahapan Hasil Pencarian Dokumen Sesuai PRISMA}
    \label{fig:prisma}
\end{figure}

Alur pencarian yang terlihat di Gambar \ref{fig:prisma} di atas, menjelaskan proses pencarian secara keseluruhan. Pada pencarian tingkat pertama, kedua basis data menghasilkan total 620 literatur. Kemudian, menggunakan pencarian tingkat 2 dengan kata kunci yang lebih tajam, dihasilkan 70 literatur yang masuk ke proses \textit{screening}. Dari 70 terdapat 12 literatur terpilih setelah melalui evaluasi judul dan abstrak. Setelah itu, setelah membaca \textit{full text} secara keseluruhan, hasil dikerucutkan lagi menjadi 5 literatur sebagai \textit{included}.


%-----------------------------------------------------------------------------%
\section{Kajian Teori}
\label{sec:2_2}
%-----------------------------------------------------------------------------%
Kajian teori pada penelitian ini disusun untuk memberikan landasan konseptual dalam menganalisis pelanggaran \textit{Service Level Agreement} (SLA) pada operasi logistik \textit{third-party logistics} (3PL), khususnya dalam konteks multi-SLA. Kajian ini mencakup konsep SLA, kinerja operasional logistik, tahapan operasional \textit{(parcel journey)}, serta pendekatan analitik berbasis data. 

Konsep-konsep tersebut digunakan untuk memahami bagaimana faktor operasional mempengaruhi keterlambatan pengiriman dan konsumsi SLA. Dalam konteks ini, \textit{machine learning} dan \textit{predictive analytics} digunakan sebagai pendekatan untuk memodelkan dan memprediksi risiko pelanggaran SLA secara kuantitatif. 

Meskipun penelitian sebelumnya telah membahas keterlambatan pengiriman dan kinerja logistik, pendekatan yang digunakan umumnya masih parsial dan belum mengintegrasikan faktor operasional multi-tahapan dalam konteks multi-SLA. Oleh karena itu, kajian teori ini menjadi dasar dalam mengidentifikasi kesenjangan penelitian dan pengembangan model prediktif pada penelitian ini.

\subsection{\textit{Service Level Agreement} (SLA) dalam Logistik}

SLA merupakan kesepakatan formal antara penyedia layanan dan pelanggan yang mendefinisikan standar kinerja layanan yang harus dipenuhi. Dalam konteks logistik, SLA umumnya mencakup indikator seperti waktu pengiriman, tingkat keberhasilan pengiriman, serta kualitas layanan yang diberikan. 

Dalam sistem logistik modern, terutama pada layanan 3PL, SLA berperan sebagai tolok ukur utama dalam menilai performa operasional. SLA tidak hanya menjadi acuan evaluasi kinerja, tetapi juga memiliki implikasi langsung terhadap aspek finansial, seperti penalti akibat keterlambatan pengiriman. 

Seiring dengan berkembangnya sistem layanan berbasis teknologi, pengelolaan SLA menjadi semakin kompleks karena melibatkan berbagai entitas dalam rantai distribusi (\textit{supply chain}), sehingga diperlukan pendekatan analitik untuk memastikan pemenuhan SLA secara konsisten \citep{engel2022}.

\subsection{\textit{Multi-Service Level Agreement} (Multi-SLA)}
Dalam ekosistem logistik e-commerce, penyedia layanan 3PL seringkali harus menangani lebih dari satu standar SLA yang berasal dari berbagai platform. Kondisi ini dikenal sebagai \textit{Multi-Service Level Agreement }(Multi-SLA).

Setiap platform \textit{e-commerce} memiliki ketentuan SLA yang berbeda, baik dalam hal waktu pengiriman, definisi keterlambatan, maupun skema penalti. Perbedaan ini menimbulkan kompleksitas dalam pengelolaan operasional karena sistem harus mampu menyesuaikan proses distribusi dengan berbagai standar layanan secara bersamaan. 

Akibatnya, pengelolaan multi-SLA membutuhkan pendekatan yang lebih adaptif dan berbasis data untuk memastikan bahwa setiap pengiriman dapat memenuhi standar SLA yang berlaku.

\subsection{Kinerja Operasional \textit{Third Party Logistics} (3PL)}
\textit{Third-party logistics} (3PL) merupakan penyedia layanan logistik yang menangani proses distribusi barang atas nama perusahaan lain. Dalam konteks \textit{e-commerce}, 3PL berperan penting dalam memastikan kelancaran aliran barang dari pengirim hingga ke pelanggan.

Kinerja operasional 3PL umumnya diukur menggunakan beberapa indikator utama, seperti:
\begin{itemize}
    \item ketepatan waktu pengiriman \textit{(on-time delivery)}
    \item tingkat keberhasilan pengiriman
    \item efisiensi operasional
\end{itemize}

Menurut \cite{book:martin2016}, efektivitas sistem logistik sangat bergantung pada kemampuan dalam mengelola aliran barang dan informasi secara terintegrasi. Dalam praktiknya, kompleksitas jaringan distribusi serta variasi permintaan menjadi tantangan utama dalam menjaga kinerja operasional.

\subsection{Tahapan Operasional dan Risiko Akumulasi Keterlambatan}
Proses distribusi dalam logistik terdiri dari beberapa tahapan operasional yang dikenal sebagai \textit{parcel journey}, yaitu \textit{first mile, middle mile,} dan \textit{last mile}.

Setiap tahapan memiliki karakteristik dan tantangan operasional yang berbeda. Keterlambatan yang terjadi pada salah satu tahapan dapat berdampak pada tahapan berikutnya, sehingga bersifat kumulatif dan membentuk akumulasi keterlambatan.

Penelitian menunjukkan bahwa tahapan \textit{last-mile} merupakan salah satu tahap yang paling kritis karena memiliki tingkat kompleksitas yang tinggi dan berpengaruh signifikan terhadap ketepatan waktu pengiriman \citep{article:boysen2021}. Dengan demikian, analisis terhadap tahapan operasional menjadi penting untuk mengidentifikasi sumber utama keterlambatan dalam proses distribusi.

\subsection{\textit{Forward} dan \textit{Return Logistics}}
Dalam sistem logistik, pengiriman dapat diklasifikasikan menjadi dua jenis utama, yaitu \textit{forward logistics} dan \textit{return logistics}. \textit{Forward logistics} mengacu pada proses pengiriman barang dari pengirim ke pelanggan, sedangkan \textit{return logistics }atau \textit{Return to Sender} (RTS) mencakup aliran balik produk dari pelanggan ke titik asal untuk pengembalian atau pemrosesan lanjutan \citep{article:prajapati2022}

Dalam rantai pasok modern, kedua aliran ini terintegrasi dalam konsep \textit{closed-loop supply chain}, yang menekankan pengelolaan simultan untuk mencapai efisiensi dan keberlanjutan. Namun, karakteristik keduanya berbeda. \textit{Forward logistics} cenderung lebih terstruktur dan dapat diprediksi, sementara return logistics memiliki tingkat ketidakpastian yang lebih tinggi dari sisi waktu, jumlah, dan kondisi produk. 

Akibatnya, \textit{return logistics} menjadi lebih kompleks karena melibatkan tahapan tambahan seperti pengumpulan kembali, inspeksi, dan pemilahan. Hal ini meningkatkan durasi proses serta potensi inefisiensi operasional. Selain itu, integrasi kedua aliran menuntut koordinasi yang baik, karena ketidakseimbangan dapat mempengaruhi kinerja sistem secara keseluruhan. Dengan demikian, perbedaan karakteristik antara \textit{forward} dan \textit{return logistics} menjadi penting dalam analisis kinerja operasional. Dalam penelitian ini, kedua alur tersebut digunakan untuk memahami variasi risiko keterlambatan berdasarkan kompleksitas tahapan operasional.

\subsection{Penerapan \textit{Machine Learning} (ML) untuk 3PL}
\textit{Machine earning} merupakan pendekatan analitik yang memungkinkan sistem untuk mempelajari pola dari data historis dan digunakan untuk melakukan prediksi. Dalam konteks logistik, ML banyak digunakan untuk menganalisis proses distribusi dan memprediksi keterlambatan pengiriman.

Pendekatan ML dapat dikategorikan menjadi \textit{supervised learning} dan \textit{unsupervised learning}, dengan metode klasifikasi sebagai salah satu pendekatan yang umum digunakan dalam prediksi kejadian seperti keterlambatan pengiriman. Penerapan ML dalam \textit{supply chain} terbukti mampu meningkatkan akurasi prediksi dan efisiensi operasional, terutama dalam menangani data yang kompleks dan bersifat non-linear \citep{article:carbonneau2008}.

\subsection{\textit{Predictive Analytics} dan Penilaian Risiko}
\textit{Predictive analytics} merupakan pendekatan analitik yang digunakan untuk memprediksi kejadian di masa depan berdasarkan data historis. Dalam penelitian ini, pendekatan ini digunakan untuk mengestimasi probabilitas terjadinya pelanggaran SLA. Penilaian risiko \textit{(risk scoring)} dilakukan dengan menghasilkan nilai probabilitas yang menunjukkan tingkat kemungkinan terjadinya keterlambatan atau SLA breach. Pendekatan ini memungkinkan pengambilan keputusan yang lebih proaktif dalam pengelolaan operasional logistik.

Menurut \cite{article:bertsimas2020}, predictive analytics memungkinkan transisi dari analisis deskriptif menuju analisis prediktif dan preskriptif, sehingga mendukung pengambilan keputusan berbasis data. Dengan demikian, kombinasi antara ML dan \textit{predictive analytics} menjadi dasar dalam pengembangan model untuk memprediksi risiko SLA breach pada penelitian ini.

\begin{landscape}

%-----------------------------------------------------------------------------%
\section{Penelitian Sebelumnya}
\label{sec:2_3}
%-----------------------------------------------------------------------------%

Penelitian sebelumnya yang relevan dengan topik prediksi risiko pelanggaran SLA pada sistem logistik telah dikaji untuk memahami pendekatan, metode, serta kontribusi yang telah dilakukan oleh peneliti terdahulu. Berdasarkan hasil SLR, diperoleh lima penelitian yang memiliki keterkaitan dengan topik penelitian ini dan akan dibahas rangkuman serta kesesuaiannya terhadap penelitian yang dilakukan.

\begin{table}[H]
\centering
\caption{Ringkasan Penelitian Sebelumnya}
\label{tab:ringkasan_penelitian}

\renewcommand{\arraystretch}{1.4}

\resizebox{\linewidth}{!}{
\begin{tabular}{|P{3cm}|P{4.5cm}|P{4.5cm}|P{4.5cm}|P{4.5cm}|P{4.5cm}|}
\hline
\textbf{Aspek} & \textbf{Penelitian 1} & \textbf{Penelitian 2} & \textbf{Penelitian 3} & \textbf{Penelitian 4} & \textbf{Penelitian 5} \\
\hline

Penulis 
& ML Delivery Cycle (2025) 
& Fuzzy Bayesian (2023) 
& Risk Network ML (2026) 
& Engel et al. (2022) 
& E-Commerce Logistics (2026) \\
\hline

Judul 
& ML-Based Delivery Cycle Classification 
& Delivery Delay Risk (Fuzzy Bayesian) 
& Risk Perception Network 
& SLA-aware AI Service Chain 
& E-Commerce Logistics Optimization \\
\hline

Tujuan 
& Klasifikasi siklus pengiriman 
& Prediksi keterlambatan 
& Analisis risiko supply chain 
& Pengelolaan SLA 
& Optimasi on-time delivery \\
\hline

Metode 
& Random Forest 
& Fuzzy Bayesian 
& Deep Learning 
& Conceptual 
& Data-driven \\
\hline

Hasil 
& Visibilitas distribusi 
& Probabilitas delay 
& Deteksi risiko 
& SLA-aware concept 
& Peningkatan performa \\
\hline

Keterbatasan 
& Tidak bahas SLA 
& Tidak SLA-based 
& Tidak spesifik 3PL 
& Tidak berbasis data 
& Tidak prediktif \\
\hline

\end{tabular}
}

\end{table}

\end{landscape}

\subsection{Perbandingan Penelitian Sebelumnya}
Membahas: 3C+2S
\begin{itemize}
    \item Compare → model yang sama
    \item Contrast → perbedaan pendekatan
    \item Criticize → kelemahan
    \item Summarize → ringkasan
    \item Synthesize → arah penelitian Anda
\end{itemize}

%-----------------------------------------------------------------------------%
\section{Kajian Metodologi}
\label{sec:2_4}
%-----------------------------------------------------------------------------%

\subsection{\textit{Predictive Analytics}}
\subsection{Model Klasifikasi}
\subsection{\textit{Logistic Regression}}
\subsection{\textit{Random Forest}}
\subsection{Evaluasi Model}
Harus ada perbandingan

%-----------------------------------------------------------------------------%
\section{Kerangka Teoretis}
\label{sec:2_5}
%-----------------------------------------------------------------------------%

\subsection{Model Konseptual Penelitian}

Model konseptual penelitian ini menjelaskan bahwa faktor operasional dalam proses distribusi, yang terdiri dari tahapan operasional (\textit{parcel journey}) seperti \textit{First Mile}, \textit{Middle Mile}, dan \textit{Last Mile}, serta dinamika \textit{reverse logistics} atau \textit{Return to Sender} (RTS), berperan sebagai variabel utama yang mempengaruhi keterlambatan dalam proses pengiriman.

Keterlambatan yang terjadi pada setiap tahapan bersifat kumulatif dan membentuk akumulasi keterlambatan. Akumulasi ini selanjutnya mempengaruhi tingkat konsumsi SLA (\textit{SLA consumption}), yaitu proporsi waktu SLA yang telah terpakai selama proses distribusi berlangsung.

Semakin tinggi tingkat konsumsi SLA, maka semakin besar probabilitas terjadinya pelanggaran SLA (\textit{SLA breach}). Oleh karena itu, hubungan antara faktor operasional, keterlambatan, dan konsumsi SLA menjadi dasar dalam memodelkan risiko pelanggaran SLA secara kuantitatif. 

\subsection{Hubungan Antar Variabel}

Hubungan antar variabel dalam penelitian ini dapat dijelaskan sebagai berikut:

\begin{enumerate}
    \item \textbf{Faktor Operasional terhadap Akumulasi Keterlambatan} \\
    Faktor operasional seperti durasi pada setiap tahapan perjalanan, jumlah tahapan (\textit{stage count}), serta status perjalanan (Forward atau RTS) mempengaruhi tingkat keterlambatan yang terjadi. Keterlambatan pada satu tahapan dapat berdampak pada tahapan berikutnya sehingga bersifat akumulatif.

    \item \textbf{Akumulasi Keterlambatan terhadap Konsumsi SLA} \\
    Akumulasi keterlambatan secara langsung meningkatkan konsumsi waktu SLA. Semakin besar keterlambatan yang terjadi, maka semakin cepat waktu SLA habis.

    \item \textbf{Konsumsi SLA terhadap Risiko Pelanggaran SLA} \\
    Konsumsi SLA yang tinggi meningkatkan probabilitas terjadinya pelanggaran SLA. Kondisi ini menjadikan konsumsi SLA sebagai indikator utama dalam memprediksi risiko keterlambatan.

    \item \textbf{Peran Jenis Perjalanan (Forward vs RTS)} \\
    Jenis perjalanan memiliki karakteristik risiko yang berbeda. RTS cenderung memiliki tingkat kompleksitas yang lebih tinggi sehingga berpotensi meningkatkan risiko pelanggaran SLA dibandingkan dengan perjalanan Forward.

    \item \textbf{Pemodelan Risiko dengan \textit{Machine Learning}} \\
    Untuk mengestimasi probabilitas pelanggaran SLA, penelitian ini menggunakan pendekatan \textit{machine learning}, khususnya metode klasifikasi seperti \textit{Logistic Regression} dan \textit{Random Forest}. Model ini memanfaatkan data historis untuk menghasilkan probabilitas risiko secara kuantitatif.
\end{enumerate}

\subsection{Output Penelitian dan Penilaian Risiko}

Berdasarkan hubungan antar variabel tersebut, penelitian ini menghasilkan mekanisme \textbf{penilaian risiko} (\textit{risk scoring}) yang digunakan untuk mengestimasi probabilitas terjadinya pelanggaran SLA pada setiap pengiriman.

Output model berupa \textit{skor risiko} yang merepresentasikan tingkat kemungkinan terjadinya SLA breach. Skor ini dapat digunakan sebagai dasar dalam sistem peringatan dini (\textit{early warning system}) untuk mendukung pengambilan keputusan operasional secara proaktif, seperti prioritisasi pengiriman dan alokasi sumber daya.

\subsection{Ilustrasi Kerangka Teoretis}

Kerangka teoretis penelitian ini secara konseptual dapat digambarkan sebagai berikut:


\begin{figure}
	\centering
	\includegraphics[width=1\textwidth]{MKSK/Tugas 1/assets/pics/kerangka_teoretis.png}
	\caption{Kerangka Teoretis Penelitian}
	\label{fig:kerangka_teori}
\end{figure}